\documentclass[a4paper]{article}

\def\npart{III}

\def\ntitle{Differential Geometry}
\def\nlecturer{A.\ Kovalev}

\def\nterm{Michaelmas}
\def\nyear{2018}

\ifx \nauthor\undefined
  \def\nauthor{Qiangru Kuang}
\else
\fi

\ifx \ntitle\undefined
  \def\ntitle{Template}
\else
\fi

\ifx \nauthoremail\undefined
  \def\nauthoremail{qk206@cam.ac.uk}
\else
\fi

\ifx \ndate\undefined
  \def\ndate{\today}
\else
\fi

\title{\ntitle}
\author{\nauthor}
\date{\ndate}

%\usepackage{microtype}
\usepackage{mathtools}
\usepackage{amsthm}
\usepackage{stmaryrd}%symbols used so far: \mapsfrom
\usepackage{empheq}
\usepackage{amssymb}
\let\mathbbalt\mathbb
\let\pitchforkold\pitchfork
\usepackage{unicode-math}
\let\mathbb\mathbbalt%reset to original \mathbb
\let\pitchfork\pitchforkold

\usepackage{imakeidx}
\makeindex[intoc]

%to address the problem that Latin modern doesn't have unicode support for setminus
%https://tex.stackexchange.com/a/55205/26707
\AtBeginDocument{\renewcommand*{\setminus}{\mathbin{\backslash}}}
\AtBeginDocument{\renewcommand*{\models}{\vDash}}%for \vDash is same size as \vdash but orginal \models is larger
\AtBeginDocument{\let\Re\relax}
\AtBeginDocument{\let\Im\relax}
\AtBeginDocument{\DeclareMathOperator{\Re}{Re}}
\AtBeginDocument{\DeclareMathOperator{\Im}{Im}}
\AtBeginDocument{\let\div\relax}
\AtBeginDocument{\DeclareMathOperator{\div}{div}}

\usepackage{tikz}
\usetikzlibrary{automata,positioning}
\usepackage{pgfplots}
%some preset styles
\pgfplotsset{compat=1.15}
\pgfplotsset{centre/.append style={axis x line=middle, axis y line=middle, xlabel={$x$}, ylabel={$y$}, axis equal}}
\usepackage{tikz-cd}
\usepackage{graphicx}
\usepackage{newunicodechar}

\usepackage{fancyhdr}

\fancypagestyle{mypagestyle}{
    \fancyhf{}
    \lhead{\emph{\nouppercase{\leftmark}}}
    \rhead{}
    \cfoot{\thepage}
}
\pagestyle{mypagestyle}

\usepackage{titlesec}
\newcommand{\sectionbreak}{\clearpage} % clear page after each section
\usepackage[perpage]{footmisc}
\usepackage{blindtext}

%\reallywidehat
%https://tex.stackexchange.com/a/101136/26707
\usepackage{scalerel,stackengine}
\stackMath
\newcommand\reallywidehat[1]{%
\savestack{\tmpbox}{\stretchto{%
  \scaleto{%
    \scalerel*[\widthof{\ensuremath{#1}}]{\kern-.6pt\bigwedge\kern-.6pt}%
    {\rule[-\textheight/2]{1ex}{\textheight}}%WIDTH-LIMITED BIG WEDGE
  }{\textheight}% 
}{0.5ex}}%
\stackon[1pt]{#1}{\tmpbox}%
}

%\usepackage{braket}
\usepackage{thmtools}%restate theorem
\usepackage{hyperref}

% https://en.wikibooks.org/wiki/LaTeX/Hyperlinks
\hypersetup{
    %bookmarks=true,
    unicode=true,
    pdftitle={\ntitle},
    pdfauthor={\nauthor},
    pdfsubject={Mathematics},
    pdfcreator={\nauthor},
    pdfproducer={\nauthor},
    pdfkeywords={math maths \ntitle},
    colorlinks=true,
    linkcolor={red!50!black},
    citecolor={blue!50!black},
    urlcolor={blue!80!black}
}

\usepackage{cleveref}



% TODO: mdframed often gives bad breaks that cause empty lines. Would like to switch to tcolorbox.
% The current workaround is to set innerbottommargin=0pt.

%\usepackage[theorems]{tcolorbox}





\usepackage[framemethod=tikz]{mdframed}
\mdfdefinestyle{leftbar}{
  %nobreak=true, %dirty hack
  linewidth=1.5pt,
  linecolor=gray,
  hidealllines=true,
  leftline=true,
  leftmargin=0pt,
  innerleftmargin=5pt,
  innerrightmargin=10pt,
  innertopmargin=-5pt,
  % innerbottommargin=5pt, % original
  innerbottommargin=0pt, % temporary hack 
}
%\newmdtheoremenv[style=leftbar]{theorem}{Theorem}[section]
%\newmdtheoremenv[style=leftbar]{proposition}[theorem]{proposition}
%\newmdtheoremenv[style=leftbar]{lemma}[theorem]{Lemma}
%\newmdtheoremenv[style=leftbar]{corollary}[theorem]{corollary}

\newtheorem{theorem}{Theorem}[section]
\newtheorem{proposition}[theorem]{Proposition}
\newtheorem{lemma}[theorem]{Lemma}
\newtheorem{corollary}[theorem]{Corollary}
\newtheorem{axiom}[theorem]{Axiom}
\newtheorem*{axiom*}{Axiom}

\surroundwithmdframed[style=leftbar]{theorem}
\surroundwithmdframed[style=leftbar]{proposition}
\surroundwithmdframed[style=leftbar]{lemma}
\surroundwithmdframed[style=leftbar]{corollary}
\surroundwithmdframed[style=leftbar]{axiom}
\surroundwithmdframed[style=leftbar]{axiom*}

\theoremstyle{definition}

\newtheorem*{definition}{Definition}
\surroundwithmdframed[style=leftbar]{definition}

\newtheorem*{slogan}{Slogan}
\newtheorem*{eg}{Example}
\newtheorem*{ex}{Exercise}
\newtheorem*{remark}{Remark}
\newtheorem*{notation}{Notation}
\newtheorem*{convention}{Convention}
\newtheorem*{assumption}{Assumption}
\newtheorem*{question}{Question}
\newtheorem*{answer}{Answer}
\newtheorem*{note}{Note}
\newtheorem*{application}{Application}

%operator macros

%basic
\DeclareMathOperator{\lcm}{lcm}

%matrix
\DeclareMathOperator{\tr}{tr}
\DeclareMathOperator{\Tr}{Tr}
\DeclareMathOperator{\adj}{adj}

%algebra
\DeclareMathOperator{\Hom}{Hom}
\DeclareMathOperator{\End}{End}
\DeclareMathOperator{\id}{id}
\DeclareMathOperator{\im}{im}
\DeclarePairedDelimiter{\generation}{\langle}{\rangle}

%groups
\DeclareMathOperator{\sym}{Sym}
\DeclareMathOperator{\sgn}{sgn}
\DeclareMathOperator{\inn}{Inn}
\DeclareMathOperator{\aut}{Aut}
\DeclareMathOperator{\GL}{GL}
\DeclareMathOperator{\SL}{SL}
\DeclareMathOperator{\PGL}{PGL}
\DeclareMathOperator{\PSL}{PSL}
\DeclareMathOperator{\SU}{SU}
\DeclareMathOperator{\UU}{U}
\DeclareMathOperator{\SO}{SO}
\DeclareMathOperator{\OO}{O}
\DeclareMathOperator{\PSU}{PSU}

%hyperbolic
\DeclareMathOperator{\sech}{sech}

%field, galois heory
\DeclareMathOperator{\ch}{ch}
\DeclareMathOperator{\gal}{Gal}
\DeclareMathOperator{\emb}{Emb}



%ceiling and floor
%https://tex.stackexchange.com/a/118217/26707
\DeclarePairedDelimiter\ceil{\lceil}{\rceil}
\DeclarePairedDelimiter\floor{\lfloor}{\rfloor}


\DeclarePairedDelimiter{\innerproduct}{\langle}{\rangle}

%\DeclarePairedDelimiterX{\norm}[1]{\lVert}{\rVert}{#1}
\DeclarePairedDelimiter{\norm}{\lVert}{\rVert}



%Dirac notation
%TODO: rewrite for variable number of arguments
\DeclarePairedDelimiterX{\braket}[2]{\langle}{\rangle}{#1 \delimsize\vert #2}
\DeclarePairedDelimiterX{\braketthree}[3]{\langle}{\rangle}{#1 \delimsize\vert #2 \delimsize\vert #3}

\DeclarePairedDelimiter{\bra}{\langle}{\rvert}
\DeclarePairedDelimiter{\ket}{\lvert}{\rangle}




%macros

%general

%divide, not divide
\newcommand*{\divides}{\mid}
\newcommand*{\ndivides}{\nmid}
%vector, i.e. mathbf
%https://tex.stackexchange.com/a/45746/26707
\newcommand*{\V}[1]{{\ensuremath{\symbf{#1}}}}
%closure
\newcommand*{\cl}[1]{\overline{#1}}
%conjugate
\newcommand*{\conj}[1]{\overline{#1}}
%set complement
\newcommand*{\stcomp}[1]{\overline{#1}}
\newcommand*{\compose}{\circ}
\newcommand*{\nto}{\nrightarrow}
\newcommand*{\p}{\partial}
%embed
\newcommand*{\embed}{\hookrightarrow}
%surjection
\newcommand*{\surj}{\twoheadrightarrow}
%power set
\newcommand*{\powerset}{\mathcal{P}}

%matrix
\newcommand*{\matrixring}{\mathcal{M}}

%groups
\newcommand*{\normal}{\trianglelefteq}
%rings
\newcommand*{\ideal}{\trianglelefteq}

%fields
\renewcommand*{\C}{{\mathbb{C}}}
\newcommand*{\R}{{\mathbb{R}}}
\newcommand*{\Q}{{\mathbb{Q}}}
\newcommand*{\Z}{{\mathbb{Z}}}
\newcommand*{\N}{{\mathbb{N}}}
\newcommand*{\F}{{\mathbb{F}}}
%not really but I think this belongs here
\newcommand*{\A}{{\mathbb{A}}}

%asymptotic
\newcommand*{\bigO}{O}
\newcommand*{\smallo}{o}

%probability
\newcommand*{\prob}{\mathbb{P}}
\newcommand*{\E}{\mathbb{E}}

%vector calculus
\newcommand*{\gradient}{\V \nabla}
\newcommand*{\divergence}{\gradient \cdot}
\newcommand*{\curl}{\gradient \cdot}

%logic
\newcommand*{\yields}{\vdash}
\newcommand*{\nyields}{\nvdash}

%differential geometry
\renewcommand*{\H}{\mathbb{H}}
\newcommand*{\transversal}{\pitchfork}
\renewcommand{\d}{\mathrm{d}} % exterior derivative

%number theory
\newcommand*{\legendre}[2]{\genfrac{(}{)}{}{}{#1}{#2}}%Legendre symbol


\DeclareMathOperator{\Gr}{Gr} % Grassmannian

\begin{document}

\begin{titlepage}
  \begin{center}
    \includegraphics[width=0.6\textwidth]{logo.jpg}\par
    \vspace{1cm}
    {\scshape\huge Mathematics Tripos \par}
    \vspace{2cm}
    {\huge Part \npart \par}
    \vspace{0.6cm}
    {\Huge \bfseries \ntitle \par}
    \vspace{1.2cm}
    {\Large\nterm, \nyear \par}
    \vspace{2cm}
    
    {\large \emph{Lectures by } \par}
    \vspace{0.2cm}
    {\Large \scshape \nlecturer}
    
    \vspace{0.5cm}
    {\large \emph{Notes by }\par}
    \vspace{0.2cm}
    {\Large \scshape \href{mailto:\nauthoremail}{\nauthor}}
 \end{center}
\end{titlepage}

\tableofcontents

\section{Manifolds}

We want to generalise curves and surfaces in \(\R^2\). A curve is a map \(\gamma: \R \supseteq I \to \R^2\) or \(\R^3\) that satisfies certain properties we'll find out in a minute. Clearly continuity is necessary but not sufficent, as evidenced by the famous Peano space-filling curve. Smoothness is not quite enough either, as \(t \mapsto (t^2, t^3)\) has a cusp at the origin. The correct requirement will be that \(\gamma\) has regular parameterisation, i.e.\ \(|\dot \gamma(t)| \neq 0\) for all \(t\).

Similarly, a surface should be defined as as a map \(r: \R^2 \supseteq D \to \R^3\) with a regular parameterisation, i.e.\ \(r \in C^\infty(D)\) and \(\left| \frac{\p r}{\p u} \times \frac{\p r}{\p v} \right| \neq 0\) for all \((u, v) \in D\).

We may follow this route and generalise to (hyper)surfaces in \(\R^n\), which will be a generalisation of classical differential geometry of curves and surfaces in \(\R^3\). The good thing is that we can readily apply calculus and it is easy to construct these objects. However, it does suffer from the prolblem of different parameterisations give rise to different geometric objects, as well as some surface requiring more than one parameterisation. Although these can be bypassed more or less, the more serious drawback is the extra technical complexity determined by the higher dimension of the ambient space.

A better concept is \emph{smooth manifolds}. We first begin with a review of topological structure, which every manifold possesses.

\begin{definition}[topological space]
  A \emph{topological space} \(M\) is a choice of class of the \emph{open sets} such that
  \begin{enumerate}
  \item \(\emptyset\) and \(M\) are open,
  \item if \(U\) and \(U'\) are open then so is \(U \cap U'\),
  \item for anly collection of open sets, the union is open.
  \end{enumerate}
\end{definition}

In this course, we always require a topological space to be Hausdorff and second countable.

\begin{definition}[local coordinate chart]\index{chart}
  A \emph{local coordinate chart} on a topological space \(M\) is a homeomorphism \(\varphi: U \to V\) where \(U \subseteq M\) and \(V \subseteq \R^d\) are open. \(U\) is a \emph{coordinate neighbourhood}.
\end{definition}

\begin{definition}[\(C^\infty\)-differentiable structure]\index{differentiable structure}\index{atlas}
  A \emph{\(C^\infty\)-differentiable structure} on a topological space \(M\) is a collection of charts \(\{\varphi_\alpha: U_\alpha \to V_\alpha\}\) where \(V_\alpha \subseteq \R^d\) for all \(\alpha\) such that
    \begin{enumerate}
    \item \(\{U_\alpha\}\) covers \(M\), i.e.\ \(M = \bigcup U_\alpha\),
    \item compatibility condition: for all \(\alpha, \beta\), \(\varphi_\beta \compose \beta_\alpha^{-1}\) is \(C^\infty\) wherever defined, i.e.\ on \(\varphi_\alpha(U_\alpha \cap U_\beta)\),
    \item maximality: if \(\varphi\) is compatible with all the \(\varphi_\alpha\)'s then \(\varphi\) is in the collection.
    \end{enumerate}
    
    The collection of charts is called an \emph{atlas}.
\end{definition}

Note that 2 implies that \(\varphi_\beta \compose \varphi_\alpha^{-1}\) is a diffeomorphism.

\begin{definition}[manifold]\index{manifold}
  A \emph{manifold} is a Hausdorff, second countable topological space with a \(C^\infty\)-differentiable structure.
\end{definition}

\begin{remark}\leavevmode
  \begin{enumerate}
  \item In practice, we almost never specify the topological structure. Instead, we can induce a topology from a \(C^\infty\) structure by declaring \(D \subseteq M\) open if and only if for all \((\varphi_\alpha, U_\alpha)\), \(\varphi_\alpha(U_\alpha \cap D)\) is open in \(\R^d\).
  \item We may replace \(C^\infty\) by \(C^k\) for \(k > 0\) finite. If we set \(k = 0\), the objects become topological manifolds. On the other hand, use \(\C^n\) and holomorphic maps we get complex manifolds.
  \item Requirement of being Hausdorff and second countable are for rather technical reasons. In some cases we may drop Hausdorffness requirement, and such examples do arise naturally. However in that case we lose uniqueness of limits. Similarly non-second countable space, such as the disjoint union of uncountably many \(\R^n\), can become manifold-like structure if we relax the definition.
  \end{enumerate}
\end{remark}

\begin{eg}\leavevmode
  \begin{enumerate}
  \item \(\R^d\) covered by the single chart \(\varphi = \id\).
  \item Unit sphere \(S^n = \{\V x = (x_0, \dots, x_n) \in \R^{n + 1}: \sum_{i = 0}^n x_i^2 = 1\}\). The charts are stereographic projections
    \begin{align}
      \varphi(\V x) &= \frac{1}{1 - x_0} (x_1, \dots, x_n) \\
      \psi(\V x) &= \frac{1}{1 + x_0} (x_1, \dots, x_n)
    \end{align}
    where \(\varphi\) is defined for all points on \(S^n\) except \((1, 0, \dots, 0)\) and similar for \(\psi\). Suppose \(\varphi(P) = u, \psi(P) = v\), then by basic geometry \(v = \psi \compose \varphi^{-1}(u) = \frac{u}{\norm{u}^2}\).
  \item Given \(M_1, M_2\) manifolds of dimension of \(d_1\) and \(d_2\), \(M_1 \times M_2\) is a manifold of dimension \(d_1 + d_2\). We define the \emph{\(n\)-torus} to be \(T_n = \underbrace{S^1 \times \dots \times S^1}_{n}\).
  \item An open subset \(U\) of a manifold \(M\) is a manifold.
  \item The \emph{real projective space},
    \[
      \R P^n = \{\text{all straight lines through \(0\) in } \R^{n + 1}\}.
    \]
    The points in the space are \(x_0 : \dots : x_n\) where \(x_i\)'s are not all zero. Note that
    \[
      x_0 : \dots : x_n = \lambda x_0 : \dots \lambda x_n
    \]
    for all \(\lambda \neq 0\). As per a previous remark, we shall induce the topology by a \(C^\infty\) structure. The charts are \((U_i, \varphi_i)\) where \(U_i = \{x_i \neq 0\}\), and
    \[
      \varphi_i(x_0 : \dots : x_n) = (\frac{x_0}{x_i}, \dots, \hat i, \dots, \frac{x_n}{x_i})
    \]
    where \(\hat i\) denotes that the \(i\)th coordinate is omitted. The \(U_i\)'s cover \(\R P^n\) so we are left to check compatibility. For \(i < j\),
    \[
      \varphi_j \compose \varphi_i^{-1} : (y_1, \dots, y_n)
      \mapsto y_1 : \dots, \underbrace{1}_{i\text{th}} : y_n
      \mapsto (\frac{y_1}{y_j}, \dots, \frac{1}{y_j}, \dots, \hat j, \dots, \frac{y_n}{y_j})
    \]
    is smooth. Since \(i\) and \(j\) are arbitrary, \(\R p^n\) is a \(n\)-manifold.

    Similarly we can check \(\C P^n\) is a \(2n\)-manifold, and \(\H P^n\) is a \(4n\)-manifold (but this is a bit tricky due to noncommutativity).
  \item Grassmannians (over \(\R\) or \(\C\)): we define \(\Gr(k, n)\) to be all \(k\)-dimensional subspaces of the \(n\)-dimensional vector space, which generalises the projective space. For real vector spaces for example, \(\R P^n = \Gr(1, n + 1)\). We can check \(\Gr(k, n)\) is a manifold of dimension \(k(n - k)\).

    The construction is a bit technical so we will give example of one chart. Let \(U\) be the \(k\)-subspaces obtainable as the span of rows of \(k \times n\) matrices of the form \((I_k \quad *)\), and the local coordinate maps a basis of such \(k\)-subspace to the \(k \times (n - k)\) block \(*\). Since the frist \(k\) rows are linearly independent, we call \(U = U_{1 < 2 < \dots < k}\). More generally the domain of charts take the form \(U_{1 \leq i_1 < \dots < i_k \leq n}\). It is an exercise to check that this gives a valid \(C^\infty\) structure.
  \item A non-example: define an equivalence relation \((x, y) \sim (\lambda x, \frac{y}{\lambda})\) for all \(\lambda \neq 0\) and define \(X := \R^2/\sim\). \(\{xy = c\}\) is one equivalence class if \(c \neq 0\). If \(c = 0\), \(\{xy = 0\}\) is three classes \(\{(x, 0): x \neq 0\}, \{(0, 0)\}, \{(0, y): y \neq 0\}\). Thus
    \[
      X \cong (-\infty, 0) \cup \{0', 0'', 0'''\} \cup (0, \infty).
    \]
    Define charts
    \[
      \varphi_i: (-\infty, 0) \cup 0^{(i)} \cup (0, \infty) \to \R
    \]
    in the obvious ways. We can check that it gives \(X\) a valid \(C^\infty\) structure, except that the induced topology is non-Hausdorff!
  \item For a non-second countable example, see example sheet 1 Q12.
  \end{enumerate}
\end{eg}

\begin{definition}[smooth map]\index{smooth map}
  Let \(M, N\) be manifolds, a continuous map \(f: M \to N\) is \emph{smooth} or \(C^\infty\) if for any \(p \in M\), there exists charts \((U, \varphi)\) on \(M\), \((V, \psi)\) on \(N\) such that \(p \in U, f(p) \in V\) and
  \[
    \psi \compose f \compose \varphi^{-1}
  \]
  is smooth wherever it is defined, i.e.\ on \(\varphi(U \cap f^{-1}(V)) \subseteq \R^n\) where \(n = \dim M\).
\end{definition}

\begin{remark}
  Note that by continuity, \(\varphi(U \cap f^{-1}(V))\) is necessarily open. Of course we can do it differently by not requiring \(f\) to be continuous a priori but instead ask the above set to be open.
\end{remark}

For any charts \(\tilde \varphi, \tilde \psi\),
\[
  \tilde \psi \compose f \compose \tilde \varphi^{-1} = (\tilde \psi \compose \psi^{-1}) \compose (\psi \compose f \compose \varphi^{-1}) \compose (\varphi \compose \tilde \varphi^{-1})
\]
is a composition of smooth map so is smooth. Thus smoothness of a map is independent of charts.

\begin{definition}[diffeomorphism]\index{diffeomorphism}
  A smooth map \(f: M \to N\) is a \emph{diffeomorphism} if \(f\) is bijective and \(f^{-1}\) is smooth. If such \(f\) exists, \(M\) and \(N\) are \emph{diffeomorphism}.
\end{definition}

\begin{remark}\leavevmode
\begin{enumerate}
\item Our definition of smoothness is a generalisation of that in calculus. More precisely, \(f: \R^n \to \R^m\) is smooth if and only if it is smooth in the calculus sense.
\item Any chart \(\varphi: U \to \R^d\) is a diffeomorphism onto its image.
\item Composition of smooth maps is smooth.
\end{enumerate}
\end{remark}

\section{Matrix Lie groups}

We can view \(\GL(n, \R)\) as an array of numbers and thus embed it in \(\R^{n^2}\). Furthermore, by considering the determinant function it is an open subset so is an \(n^2\)-manifold. Matrix multiplication is obviously smooth. In the same vein we have \(\GL(n, \C)\), \(\SL(n, \R)\) etc as mannifolds with smooth multiplications.

\begin{definition}[Lie group]\index{Lie group}
  A group \(G\) is a \emph{Lie group} if \(G\) is a manifold and has a compatible group structure, i.e.\ the map \((\sigma, \tau) \mapsto \sigma\tau^{-1}\) is smooth.
\end{definition}

Before Lie groups, let's have a short digression in analysis to discuss the exponential map. For a complex \(n \times n\) matrix \(A = (a_{ij})\), define a norm
\[
  |A| = n \cdot \max_{ij} |a_{ij}|
\]
where the \(n\) in front is such that \(|AB| \leq |A| \cdot |B|\). We now define
\[
  \exp(A) := I + A + \frac{1}{2} A^2 + \dots + \frac{1}{n!} A^n + \dots
\]

Of course we have to check the series makes sense: in fact \(\exp A\) is absolutely convergent for all \(A\) as \(\left| \frac{A^n}{n!} \right| \leq \frac{|A|^n}{n!}\). The series is also uniformly convergent on any compact set, by Weierstrass \(M\)-test. Therefore \(\exp\) is a well-defined continuous map.

In fact, the map is smooth although the proof is technical. A sketch of the proof is like this: note that \(f: A \to A^m\) is smooth for all \(m \in \N\). For \(m = 2\), \(df_A: H \mapsto HA + AH\) so \(\norm{df_A} \leq 2 |A|\) so for all \(m\),
\[
  \norm{df_A} \leq m |A|^{m - 1}
\]

We can thus term-by-term differentiate \(\exp\) and get a locally uniformly convergent series and use the above estimate to bound derivatives.












\printindex
\end{document}

% https://www.dpmms.cam.ac.uk/~agk22/teaching.html